\section{Introduction}
\label{sec:intro}

Modern LLM agents re-read long, mostly-static instructions on every turn---a system policy, a tool
specification, retrieved documents. Key/value (KV) caching makes this affordable by reusing the prefill
across turns, but only across an \emph{exact} shared prefix. The moment one token changes inside the reused
region---a timestamp, a user id, an order's status---the keys and values of \emph{every} later token are
invalidated, because each attended to the token that changed. The de-facto workaround---\emph{hoisting} all
mutable content to the end so the static prefix stays cache-aligned---pushes an inference-layer constraint
into the application layer: fields referenced in several places, nested sub-agent prompts, and dynamically
assembled contexts cannot all be cleanly hoisted, and the application must enumerate every mutable field in
advance.

This paper starts from a concrete puzzle. The region of the cache \emph{before} a field is, by
construction, independent of the field's value---we measure a key/value deviation of exactly $0.0$ when
the field changes. One might therefore hope to \emph{surgically} refresh only the field's own keys and
values, leave the rest of the cache stale, and pay almost nothing. We find this fails completely: the
model's decision reverts to the \emph{old} field value, as if the edit never happened
(\Cref{fig:teaser}b).

\paragraph{The discovery.} The reason, which we establish causally, is that transformers do not defer
their reasoning to decode time. At \emph{prefill} the model already computes the
\emph{field-conditioned conclusion} and writes it onto downstream tokens---disproportionately onto
aggregator/delimiter tokens that later positions attend through. The decision then reads these
\emph{notes}, not the field itself: across models the field's own KV causally drives less than $1\%$ of
the decision, while the downstream notes drive essentially all of it. The KV cache is best understood not
as a frozen byproduct of prefill but as a \emph{notebook of memoized conclusions} (\Cref{fig:teaser}a).

\paragraph{Two capabilities from one mechanism.} This reframing is generative. (1) If the conclusion is
already written downstream, then \emph{editing} a field means amending the notes, not recomputing them. The
cheap and robust fix is a one-line salient \emph{erratum} that overrides the stale notes; one can instead
recompute the affected downstream notes (reliably for the full affected suffix, cheaply-but-unreliably for
a top-$K$, \fieldsel{}@$K$). An even cheaper in-place refresh of the field alone is gated by
\emph{chain-of-thought}: with a reasoning chain the model re-reads the field and the edit suffices, without
it the edit is ignored. (2) If the notes are localized and position-portable, then a reusable skill can be
\emph{composed} into a new context by repositioning and splicing its cached notes---no recompute. A
\emph{keystone} experiment, editing a field \emph{inside} a transplanted skill, shows the two operations
act on a single substrate. \Cref{fig:overview} previews both capabilities.

\paragraph{Contributions.}
\begin{itemize}[leftmargin=1.2em,itemsep=1pt,topsep=2pt]
\item \textbf{A mechanism} (\Cref{sec:mechanism})---\emph{attention-mediated memoized inference}---established
by four causal probes (locality patching, suffix concentration, linear probing, circuit knockout) and four
further controls (content/conclusion dissociation, layer-timing, specificity, note-injection;
\Cref{app:deepmech}), replicated across \emph{four model families} (Qwen3, Llama-3.1, Gemma-2, Mistral),
resolved to a \emph{component-level circuit}---named read/write heads, a causal conclusion direction, an SAE
feature, attention-vs-MLP, and causal scrubbing (\Cref{app:circuit})---and connected to the delimiter-token
aggregation seen in interpretability.
\item \textbf{An editing capability} (\Cref{sec:editable}): naive KV editing fails; the erratum /
\texttt{field+erratum} fix matches the hoist-to-end oracle without prompt surgery; an analysis of when
the ${\sim}1\%$-compute in-place edit suffices (it requires reasoning and is strongly
model-dependent); and a head-to-head with \emph{weight}
editing (ROME, LoRA) showing it is the wrong tool for mutable per-request state (global contamination,
collateral, $30$--$50\times$ slower).
\item \textbf{A composing capability} (\Cref{sec:composable}): position-portable transplant of precompiled
skills, $O(L)$ vs.\ $O(L^2)$ TTFT ($13.9\times$ at $32$k), with a seam-repair knob---honestly credited to
prior caching work, with our addition being the mechanism that explains it and a correctness lens.
\item \textbf{The unification} (\Cref{sec:keystone}): the keystone and a unified edit+compose agent over
thirteen models.
\item \textbf{An application: user memory} (\Cref{sec:memory}): the large, mutable user-memory
document is both composed (precompiled, repositioned, spliced) and edited (in-place / erratum)
as one substrate---decision-faithful to full recompute at $2.3$--$4.3\times$ lower
time-to-first-token, validated to $70$B and on real long-conversation memory (LoCoMo, transplant
$\equiv$ full recompute in QA accuracy)---with a pre-registered, statistically-controlled evaluation.
\item \textbf{Reach and scope} (\Cref{sec:reach}): validation across scale, MoE, quantization, and
multimodal image caches; small adapters for MLA and interleaved M-RoPE; a fix for sliding-window; and an
analysis of where the substrate holds up to the 2026 sparse/compressed-attention frontier.
\item \textbf{Systems payoff} (\Cref{sec:systems}): a real agentic environment and a comprehensive online
vLLM serving benchmark (V1 engine, continuous batching, Poisson load)---$98.5\%$ vs.\ $1\%$ prefix-cache
hit-rate, $53$--$398\times$ lower $p90$ TTFT, and throughput gains that grow with load to $14.5\times$.
\end{itemize}

\paragraph{Roadmap.} \Cref{sec:mechanism} establishes the mechanism;
\Cref{sec:editable,sec:composable,sec:keystone} derive editing, composing, and their unification;
\Cref{sec:memory} applies both to user memory; and \Cref{sec:reach,sec:systems} map the substrate's reach
and the end-to-end systems payoff.

\begin{figure}[t]
\centering
\includegraphics[width=\linewidth]{fig1_teaser}
\caption{\textbf{Models take notes at prefill.} (a) At prefill the model memoizes the field-conditioned
conclusion onto downstream aggregator tokens (orange); at decode the decision reads those notes (blue).
(b) Consequently, surgically editing the field's own KV is ignored (without a reasoning chain), but the
decision is recovered cheaply: recompute the \emph{affected} downstream notes, or---cheaper and robust---
append a salient \emph{erratum} (P(new decision), no-CoT model; with chain-of-thought the field-only
refresh alone also works, \Cref{fig:overview}c). (c) Because the notes are position-portable, a precompiled
skill can be pasted into a new context in $O(L)$ rather than $O(L^2)$ time.}
\label{fig:teaser}
\end{figure}

\begin{figure}[t]
\centering
\includegraphics[width=0.92\linewidth]{fig0_overview}
\caption{\textbf{Editing and composing, previewed} (Qwen3-8B unless noted; detail in
\Cref{sec:mechanism,sec:editable,sec:keystone}). (a) \emph{KV editing landscape}: naive edits (stale,
field-only without CoT, CacheBlend) fail, while the append-only \texttt{erratum}/\texttt{field+erratum}
reach full-reprefill correctness cheaply and robustly (hoisting also works but needs prompt surgery).
(b) \emph{Recompute the affected notes}: recovery climbs as more of the post-field \emph{affected suffix} is
recomputed---reliable in the limit, but the cheap top-$K$ version (\texttt{field+selective}@$K$) is
unreliable (\Cref{fig:edit}). (c) \emph{Chain-of-thought, not model size, gates the cheapest edit}: with CoT
the near-free field-only refresh alone recovers the decision ($1.00$); without CoT it is ignored ($0.00$);
the erratum fixes both. (d) \emph{Composing $+$ editing is lossless}: the unified agent stays
decision-identical to full recompute across models.}
\label{fig:overview}
\end{figure}
